%% hopfion-framework/quantum/schrodinger.tex
%% Source: main_paper13.tex
%% Status: published
%% Last updated: 2026-07-05
%% Condensate-modified Hamiltonian, chameleon screening, energy corrections,
%% free-electron mass anisotropy, and the Schrodinger equation from
%% condensate quantisation (translation fibration, exact BO separation).
%% Depends: quantum/hilbert_space.tex (thm:quantisation, lem:momentum)


%════════════════════════════════════════════════════════════════════
% The Condensate-Modified Hydrogen Hamiltonian — Paper XIII, Sec. 2-3
%════════════════════════════════════════════════════════════════════

\begin{remark}[Radial dependence]
\label{P13:rem:radial_dep}
Equation~\eqref{P13:eq:rho_atom} gives $\bst\rho_{\rm atom}(r)\propto r^{-4}$:
the condensate density grows rapidly toward the nucleus.
At $r = 5a_0$ (outer edge of the $n=2$ shell):
$\bst\rho_{\rm atom}\approx4.1\times10^{22}$, still $\gg1$.
The chameleon limit is an excellent approximation throughout the atom.
\status{published}
\source{P13:rem:radial_dep}
\end{remark}
\status{published}
\source{P13:rem:radial_dep}

\begin{theorem}[Condensate-modified hydrogen Hamiltonian]
\label{P13:thm:modified_H}
The effective Hamiltonian for the electron gradient-knot in the
condensate background near a proton is
\begin{equation}
  \boxed{
  \Heff
  \;=\; -\frac{\hbar^2}{2m_e}\nabla^2 - \frac{e^2}{r}
  \;-\; \frac{\hbar^2}{2m_e}\,\frac{\epsr}{r^2}\,L^2,
  }
  \label{P13:eq:H_eff}
\end{equation}
where $\epsr = 1/[\vp^6(1+\bst\rho(r))]$ and $\rho(r)$ is the
condensate density at radius $r$.
The condensate correction vanishes identically for radial motion
($\Seff(0,\rho)=0$ by $\sin^4\!0=0$) and acts only on angular
kinetic energy.
\status{published}
\source{P13:thm:modified_H}
\end{theorem}
\status{published}
\source{P13:thm:modified_H}

\begin{remark}[Zero radial correction]
\label{P13:rem:zero_radial}
The exact vanishing of the condensate correction for radial motion
is a structural consequence of the Hopf map Jacobian:
$\Seff(\theta)\propto\sin^4\!\theta$, and $\sin\!0=0$ exactly.
The condensate does not resist radial propagation at all.
This is not an approximation --- it is built into the geometry
of the Hopf fibration.
\status{published}
\source{P13:rem:zero_radial}
\end{remark}
\status{published}
\source{P13:rem:zero_radial}

%════════════════════════════════════════════════════════════════════
% Chameleon Screening and the Origin of Spherical Harmonics — Sec. 4
%════════════════════════════════════════════════════════════════════
% thm:chameleon_screening is filed canonically in quantum/chameleon.tex
% (see that file for the full statement); referenced here by \ref only.

\begin{remark}[The same mechanism as the cosmological constant]
\label{P13:rem:same_mechanism}
The cosmological constant hierarchy is resolved in Paper~VII
by the chameleon mechanism: the condensate vacuum energy is present in full
but screened from gravitational coupling at cosmological density by
the denominator $1/(1+\bst\rho_\infty)$.
Theorem~\ref{P13:thm:chameleon_screening} (quantum/chameleon.tex) uses the same
denominator to derive spherical harmonics inside atoms.
The two results --- the cosmological constant and orbital shapes ---
share a single algebraic origin.
\status{published}
\source{P13:rem:same_mechanism}
\depends{quantum/chameleon.tex (P13:thm:chameleon_screening)}
\end{remark}
\status{published}
\source{P13:rem:same_mechanism}

\begin{remark}[Quantitative screening ratio]
\label{P13:rem:screening_ratio}
The ratio of condensate density inside the atom to the cosmic attractor:
\[
  \frac{\bst\rhoa(a_0)}{\bst\rhoinf}
  \;=\; \frac{\aem^6 m_e^4}{2\Lcond^4\,\vp}
  \;\approx\; \frac{2.57\times10^{25}}{1.618}
  \;\approx\; 1.6\times10^{25}.
\]
The condensate is $1.6\times10^{25}$ times denser inside the atom
than in the cosmic vacuum.
The suppression of the condensate correction inside atoms relative
to the cosmic background is the same factor:
$\eps(a_0)/\eps_\infty \approx 10^{-27}/0.021 \approx 10^{-25}$.
\status{published}
\source{P13:rem:screening_ratio}
\end{remark}
\status{published}
\source{P13:rem:screening_ratio}

\begin{proposition}[Universal chameleon screening in multi-electron atoms]
\label{P13:prop:multi_electron_screening}
For any atom with $Z\ge1$ electrons, the chameleon screening condition
$\bst\rho_{\rm atom}^{(Z)}(r)\gg1$ holds for all $r$ within the
electronic density.
Consequently:
\begin{enumerate}[noitemsep,label=(\roman*)]
\item The orbital shapes are exactly the standard spherical harmonics
  $Y_\ell^m(\theta,\varphi)$ for every electron in every atom.
\item The effective Hamiltonian for the $i$-th electron reduces to
  the standard Hartree--Fock operator:
  \begin{equation}
    \Heff^{(i)} \;\to\; -\frac{\hbar^2}{2m_e}\nabla^2
    - \frac{Z_{\rm eff}(r)\,e^2}{r}
    + V_{\rm HF}^{(i)},
    \label{P13:eq:H_HF_limit}
  \end{equation}
  where $V_{\rm HF}^{(i)}$ includes the Hartree and exchange potentials.
\end{enumerate}
\status{published}
\source{P13:prop:multi_electron_screening}
\end{proposition}
\status{published}
\source{P13:prop:multi_electron_screening}

\begin{proposition}[Condensate energy corrections in multi-electron atoms]
\label{P13:prop:multi_electron_energy}
In the chameleon limit, the condensate
correction to the energy of orbital $(n,\ell)$ in a $Z$-electron atom is
\begin{equation}
  \Delta E_{n\ell}^{(Z)}
  \;=\; -\frac{\hbar^2\,\ell(\ell+1)}{2m_e}
        \left\langle\frac{\eps^{(Z)}(r)}{r^2}\right\rangle_{n\ell},
  \label{P13:eq:delta_E_Z}
\end{equation}
where $\eps^{(Z)}(r)=1/[\vp^6(1+\bst\rho_{\rm atom}^{(Z)}(r))]$.
Two regimes apply:
\begin{enumerate}[noitemsep,label=(\roman*)]
\item \textbf{Inner electrons} ($Z_{\rm eff}\approx Z$, $r\sim a_0/Z$):
  $\Delta E_{n\ell}^{(Z,\,\rm inner)}\approx Z^{-4}\,\Delta E_{n\ell}^{(Z=1)}$,
  since $\eps^{(Z)}\approx\eps^{(1)}/Z^2$ and
  $\langle r^2\rangle_{n\ell}^{(Z)}\approx\langle r^2\rangle_{n\ell}^{(1)}/Z^2$.
  The correction is \emph{more suppressed} for heavier atoms.
\item \textbf{Outer (valence) electrons} ($Z_{\rm eff}\approx 1$, $r\sim n^2 a_0$):
  since $Z_{\rm eff}\to1$ in the outer shells,
  $\bigl|\Delta E_{n\ell}^{(Z,\,\rm outer)}\bigr|\lesssim\bigl|\Delta E_{n\ell}^{(Z=1)}\bigr|$.
\end{enumerate}
In both cases, the universal upper bound is
\begin{equation}
  \bigl|\Delta E_{n\ell}^{(Z)}\bigr|
  \;\lesssim\; \bigl|\Delta E_{n\ell}^{(Z=1)}\bigr|
  \;\sim\; 10^{-24}\,\mathrm{eV}\times\ell(\ell+1),
  \label{P13:eq:delta_E_bound}
\end{equation}
suppressed $\sim10^{20}$ below the fine structure for all atoms.
\status{published}
\source{P13:prop:multi_electron_energy}
\end{proposition}
\status{published}
\source{P13:prop:multi_electron_energy}

\begin{corollary}[The Hartree--Fock equations are exact]
\label{P13:cor:HF_exact}
Proposition~\ref{P13:prop:multi_electron_screening} establishes that the
condensate correction to the Hartree--Fock Hamiltonian is
$O(\eps^{(Z)})\lesssim O(10^{-27})$, negligible at any measured precision.
Standard quantum chemistry methods --- Hartree--Fock, DFT, coupled cluster,
configuration interaction --- are therefore all derivable within the
condensate framework, with no condensate correction needed.
In particular, the Aufbau filling order and the structure of the
periodic table (addressed via the $\vp$-spiral in Paper~XI)
are unaffected by the condensate.
\status{published}
\source{P13:cor:HF_exact}
\end{corollary}
\status{published}
\source{P13:cor:HF_exact}

\begin{theorem}[Atomic physics from the condensate]
\label{P13:thm:atomic_foundation}
This paper, taken together, derives from the condensate framework:
\begin{enumerate}[noitemsep]
\item The modified Hamiltonian with a single condensate correction
  (Theorem~\ref{P13:thm:modified_H}).
\item Spherical harmonics as exact orbital shapes for all atoms and all $Z$
  (Theorem~\ref{P13:thm:chameleon_screening} and
  Proposition~\ref{P13:prop:multi_electron_screening}).
\item Condensate energy corrections bounded by $\sim10^{-24}\,\mathrm{eV}$,
  undetectable in any atom (Propositions~\ref{P13:prop:energy_correction}
  and~\ref{P13:prop:multi_electron_energy}).
\item The Schr\"odinger equation from condensate quantisation,
  exact to $O(\eps_{\rm cond})\approx10^{-27}$
  (Theorem~\ref{P13:thm:schrodinger}, Corollary~\ref{P13:cor:improved_BO}).
\item Pauli exclusion and spin-$\tfrac{1}{2}$ from $Q=2$ topology
  (Theorems~\ref{P13:thm:spin_half}--\ref{P13:cor:pauli_exclusion},
  quantum/spin\_statistics.tex).
\item The spin assignments $J=Q/4$ for all Standard Model particles
  (Theorem~\ref{P13:thm:spin_tower}, quantum/spin\_statistics.tex).
\end{enumerate}
Standard atomic and molecular quantum mechanics --- wavefunctions,
selection rules, the periodic table, all of quantum chemistry ---
follows from the condensate framework without additional postulates.
\status{published}
\source{P13:thm:atomic_foundation}
\depends{quantum/spin_statistics.tex}
\end{theorem}
\status{published}
\source{P13:thm:atomic_foundation}

%════════════════════════════════════════════════════════════════════
% Condensate Corrections to Hydrogen Energy Levels — Sec. 5
%════════════════════════════════════════════════════════════════════

\begin{proposition}[Condensate energy correction to hydrogen levels]
\label{P13:prop:energy_correction}
In first-order perturbation theory with $H' = \Hcond$,
the condensate correction to the hydrogen energy level $|n,\ell\rangle$ is
\begin{equation}
  \Delta E_{n\ell}
  \;=\; \left\langle n,\ell\left|H'\right|n,\ell\right\rangle
  \;=\; -\frac{\hbar^2\ell(\ell+1)}{2m_e}\,
        \left\langle\frac{\epsr}{r^2}\right\rangle_{n\ell}.
  \label{P13:eq:delta_E_general}
\end{equation}
In the chameleon limit $\bst\rhoa\gg1$, inserting~\eqref{P13:eq:eps_chameleon}:
\begin{equation}
  \Delta E_{n\ell}
  \;\approx\; -\frac{\ell(\ell+1)\,\Lcond^4}{\vp^6\,\aem^4\,m_e^3}
              \cdot\frac{\langle r^2\rangle_{n\ell}}{a_0^4},
\end{equation}
where $\langle r^2\rangle_{n\ell} = (n^2/2)(5n^2+1-3\ell(\ell+1))\,a_0^2$
is the exact hydrogen expectation value.
For the $2p$ state ($n=2$, $\ell=1$):
\begin{equation}
  \boxed{\Delta E_{2p}
  \;\approx\; -\frac{2\,\Lcond^4\cdot30}{\vp^6\,\aem^4\,m_e^3}
  \;\approx\; -1.77\times10^{-24}\,\mathrm{eV}.}
  \label{P13:eq:delta_E_2p}
\end{equation}
This is $\sim1.95\times10^{-20}$ times the fine-structure
splitting of the $2p$ level
($\Delta E_{\rm fs}\sim\aem^2\ERy/n^3\approx9.1\times10^{-5}\,\mathrm{eV}$).
\status{published}
\source{P13:prop:energy_correction}
\end{proposition}
\status{published}
\source{P13:prop:energy_correction}

\begin{remark}[Structural form of the correction]
\label{P13:rem:structure}
The condensate correction~\eqref{P13:eq:delta_E_general}:
\begin{enumerate}[noitemsep]
\item Is zero for $\ell=0$ states (s-orbitals), exactly:
  $L^2|n,0\rangle = 0$.
\item Breaks the $\ell$-degeneracy of hydrogen energy levels
  (the SO(4) accidental degeneracy): states with the same $n$
  but different $\ell$ receive different corrections.
  This is a condensate analogue of the Lamb shift, at $\sim10^{-20}$
  smaller magnitude.
\item Scales as $\Lcond^4/m_e^3 \propto \TCMB^4/m_e^3$:
  it would grow as the CMB temperature increases, meaning the
  condensate correction was larger in the early universe.
\end{enumerate}
\status{published}
\source{P13:rem:structure}
\end{remark}
\status{published}
\source{P13:rem:structure}

%════════════════════════════════════════════════════════════════════
% Free-Electron Mass Anisotropy in the Cosmic Background — Sec. 6
%════════════════════════════════════════════════════════════════════

\begin{proposition}[Free-electron mass anisotropy]
\label{P13:prop:mass_anisotropy}
For a free electron propagating in the cosmic condensate background
at the attractor density $\rhoinf$, the effective mass depends on
the direction of propagation relative to the local condensate gradient:
\begin{align}
  m_{\rm radial}
  &= m_e
  &\text{(exactly, for propagation parallel to } \nabla\rho\text{)},
  \label{P13:eq:m_radial}
  \\
  m_{\rm tangential}
  &= m_e\left(1 + \frac{1}{\vp^8}\right)
   = m_e\,\frac{\vp^8+1}{\vp^8}
  &\text{(perpendicular to } \nabla\rho\text{)}.
  \label{P13:eq:m_tangential}
\end{align}
The fractional mass anisotropy is
\begin{equation}
  \boxed{\frac{\Delta m}{m_e}
  \;=\; \frac{m_{\rm tangential} - m_{\rm radial}}{m_e}
  \;=\; \frac{1}{\vp^8}
  \;\approx\; 2.129\%.}
  \label{P13:eq:mass_anisotropy}
\end{equation}
The radial-to-tangential effective mass ratio is
\begin{equation}
  \frac{m_{\rm radial}}{m_{\rm tangential}}
  \;=\; \frac{\vp^8}{\vp^8+1}
  \;\approx\; 0.97916,
  \label{P13:eq:mass_ratio}
\end{equation}
exact in terms of $\vp$ with no free parameters.
\status{prediction}
\source{P13:prop:mass_anisotropy}
\end{proposition}
\status{prediction}
\source{P13:prop:mass_anisotropy}

\begin{remark}[Exact form and approximations]
\label{P13:rem:exact_vs_approx}
The result~\eqref{P13:eq:mass_anisotropy} uses two exact golden-ratio
identities: $\bst\rhoinf=\vp$ (the condensate self-consistency
condition, Paper~I) and $1+\vp=\vp^2$
(the pentagon identity).
The mass ratio $\vp^8/(\vp^8+1)$ involves no free parameters.
The $2.129\%$ anisotropy follows from these alone.
\status{published}
\source{P13:rem:exact_vs_approx}
\end{remark}
\status{published}
\source{P13:rem:exact_vs_approx}

\begin{remark}[Direction of the preferred axis and experimental access]
\label{P13:rem:experimental_access}
The condensate gradient $\nabla\rho$ near a gravitational source
(Earth, Sun, Galaxy) is radially directed.
In empty space far from any mass, the large-scale condensate
gradient is determined by the local matter distribution and the
CMB dipole.
A free electron (or any lepton) propagating radially versus
tangentially relative to the galactic centre, or relative to
the CMB dipole direction, should show a $2.129\%$ difference
in effective mass.
The natural experimental probe is a precision Penning trap:
the cyclotron frequency $\omega_c = eB/m_{\rm eff}$ would depend
on the orientation of the trap relative to the local condensate
gradient at the $2.129\%/\vp^8$ level.

Current Penning-trap measurements of the electron $g$-factor
(Hanneke et al.\ 2008) have fractional precision
$\sim10^{-13}$, which could in principle detect a $2\%$ effect
if the anisotropy is correctly modelled.
A dedicated anisotropy search --- varying trap orientation
over a year to scan the local condensate gradient direction ---
would provide a clean test.
\status{prediction}
\source{P13:rem:experimental_access}
\end{remark}
\status{prediction}
\source{P13:rem:experimental_access}

\begin{remark}[Relation to the $\alpha$--$m_e$ tracking relation]
\label{P13:rem:tracking_relation}
Paper~IV establishes the tracking relation
$\delta m_e/m_e = -\delta\aem/(\pi\aem) \approx -43.6\,\delta\aem/\aem$:
as $\aem$ is refined, $m_e$ tracks.
The mass anisotropy~\eqref{P13:eq:mass_anisotropy} is a distinct
effect: the anisotropy is in the effective mass during propagation,
not in the rest mass $m_e$.
The rest mass is set by the tower level and WZW T-phase (Papers~I, IV)
and is isotropic.
The propagation effective mass acquires the $1/\vp^8$ anisotropy
from the condensate medium.
\status{published}
\source{P13:rem:tracking_relation}
\end{remark}
\status{published}
\source{P13:rem:tracking_relation}

%════════════════════════════════════════════════════════════════════
% Schrodinger Equation from Condensate Quantisation — Sec. 7
%════════════════════════════════════════════════════════════════════

\begin{definition}[One-particle states]
\label{P13:def:one_particle}
A one-particle state is a state of the form
\begin{equation}
  |\Psi_e\rangle \;=\; \psi(\mathbf{R})\otimes|0_{\rm internal}\rangle,
  \label{P13:eq:one_particle_state}
\end{equation}
where $\psi\in L^2(\mathbb{R}^3)$ is the centre-of-mass wavefunction
and $|0_{\rm internal}\rangle\in\mathcal{H}_{\rm internal}$
is the ground state of the condensate knot's internal structure.
\status{published}
\source{P13:def:one_particle}
\end{definition}
\status{published}
\source{P13:def:one_particle}

\begin{lemma}[Momentum operator on the one-particle sector]
\label{P13:lem:momentum_operator}
On the one-particle sector~\eqref{P13:eq:one_particle_state},
the operator corresponding to the Noether momentum~\eqref{P13:eq:total_momentum}
acts as
\begin{equation}
  \hat{P}_i\,\psi(\mathbf{R})\otimes|0_{\rm internal}\rangle
  \;=\; \bigl(-i\hbar\,\partial/\partial R_i\,\psi(\mathbf{R})\bigr)
        \otimes|0_{\rm internal}\rangle.
  \label{P13:eq:momentum_on_sector}
\end{equation}
\status{published}
\source{P13:lem:momentum_operator}
\depends{quantum/hilbert_space.tex (P10:lem:momentum)}
\end{lemma}
\status{published}
\source{P13:lem:momentum_operator}

\begin{theorem}[Born--Oppenheimer separation for the condensate]
\label{P13:thm:born_oppenheimer}
In the adiabatic limit $\eps_{\rm BO}\to0$, the condensate Hamiltonian~\eqref{P13:eq:H_decomp}
restricted to one-particle states~\eqref{P13:eq:one_particle_state} with
$|0_{\rm internal}\rangle$ in the ground state of $\hat{H}_{\rm internal}$ becomes
\begin{equation}
  \hat{H}_{\rm eff}
  \;=\; \frac{\hat{\mathbf{P}}^2}{2m_e}
  \;+\; \hat{V}(\mathbf{R})
  \;+\; O\!\left(\eps_{\rm BO}^2\right),
  \label{P13:eq:Heff_BO}
\end{equation}
where $m_e = \Lcond\,\vp^{20}\,e^{4\pi-3/400-\aem/\pi}$
is the condensate-derived electron mass (Papers~I,~IV).
\status{published}
\source{P13:thm:born_oppenheimer}
\end{theorem}
\status{published}
\source{P13:thm:born_oppenheimer}

\begin{theorem}[Schr\"odinger equation from condensate quantisation]
\label{P13:thm:schrodinger}
In the one-particle sector of the condensate Hilbert space,
the time evolution of the electron's centre-of-mass wavefunction
$\psi(\mathbf{R},t)\in L^2(\mathbb{R}^3)$ satisfies
\begin{equation}
  i\hbar\,\frac{\partial\psi}{\partial t}
  \;=\;
  \left[-\frac{\hbar^2}{2m_e}\nabla^2 + V(\mathbf{R})\right]\psi
  \;+\; O\!\left(\eps_{\rm cond}\right),
  \label{P13:eq:schrodinger_derived}
\end{equation}
where $m_e$ is the condensate-derived electron mass,
$V(\mathbf{R})$ is any external potential,
and $\eps_{\rm cond}\sim(\Lcond/m_e\aem)^4\sim10^{-27}$ is the
condensate correction from Theorem~\ref{P13:thm:modified_H}.
The $O(\eps_{\rm BO}^2)$ Born--Oppenheimer remainder vanishes exactly
by the translation fibration (Theorem~\ref{P13:thm:NA_vanishing}
and Corollary~\ref{P13:cor:improved_BO}); the error bound
$O(\eps_{\rm cond})\approx O(10^{-27})$ is a factor of
$\sim10^{17}$ tighter than the formal BO bound.
\status{published}
\source{P13:thm:schrodinger}
\depends{quantum/hilbert_space.tex (P10:thm:quantisation)}
\end{theorem}
\status{published}
\source{P13:thm:schrodinger}

\begin{remark}[What the derivation establishes]
\label{P13:rem:what_derived}
Theorem~\ref{P13:thm:schrodinger} shows that the Schr\"odinger equation
is not an additional postulate of the condensate framework.
It is a consequence of four structural results already present:
\begin{enumerate}[noitemsep]
  \item The Hilbert space $\mathcal{H}=L^2(\mathcal{C}_Q^{\rm red},d\mu)$
    and time-evolution equation (Paper~X, Theorem thm:quantisation).
  \item The translation fibration $\mathcal{C}_Q^{\rm red}=\mathbb{R}^3\times
    \mathcal{C}_{\rm internal}$ (equation~\eqref{P13:eq:fibration}).
  \item The Noether momentum $\hat{P}_i=-i\hbar\partial/\partial R_i$
    on the one-particle sector (Lemma~\ref{P13:lem:momentum_operator}).
  \item The Born--Oppenheimer adiabatic separation at
    $\eps_{\rm BO}=\aem^2/(2\vp)\approx1.6\times10^{-5}$
    (Theorem~\ref{P13:thm:born_oppenheimer}).
\end{enumerate}
The electron mass $m_e$ appearing in~\eqref{P13:eq:schrodinger_derived}
is not an input to this derivation: it is the condensate mass
from Papers~I and~IV, derived from $\TCMB$ alone.
\status{published}
\source{P13:rem:what_derived}
\end{remark}
\status{published}
\source{P13:rem:what_derived}

\begin{remark}[Adiabatic correction estimate --- initial bound]
\label{P13:rem:BO_correction}
Theorem~\ref{P13:thm:born_oppenheimer} states the Born--Oppenheimer error as
$O(\eps_{\rm BO}^2)$ where $\eps_{\rm BO}=\aem^2/(2\vp)\approx1.6\times10^{-5}$,
giving a formal bound of $\eps_{\rm BO}^2\approx2.7\times10^{-10}$
on the fractional energy correction.
This bound is sharp for a generic BO problem (e.g.\ a diatomic molecule),
where the derivative coupling $d_{0k}=\langle k|\partial/\partial R|0\rangle\neq0$.
We now show that in the condensate framework this coupling vanishes
\emph{identically}, improving the bound by a factor of $\sim10^{17}$.
\status{published}
\source{P13:rem:BO_correction}
\end{remark}
\status{published}
\source{P13:rem:BO_correction}

\begin{theorem}[Exact vanishing of non-adiabatic coupling]
\label{P13:thm:NA_vanishing}
The Born--Oppenheimer derivative coupling between any two internal states
of the $Q=2$ gradient-knot vanishes identically:
\begin{equation}
  \boxed{d_{0k} \;\equiv\; \langle k_{\rm internal}|\,\partial/\partial R_i\,|0_{\rm internal}\rangle
  \;=\; 0
  \qquad \text{for all } k \text{ and all } i=1,2,3.}
  \label{P13:eq:d0k_zero}
\end{equation}
Consequently, the diagonal Born--Oppenheimer correction (DBOC) also vanishes:
\begin{equation}
  \mathrm{DBOC}
  \;\equiv\; -\frac{\hbar^2}{2m_e}\sum_{k\neq0}\frac{|d_{0k}|^2}{E_k-E_0}
  \;=\; 0
  \label{P13:eq:DBOC_zero}
\end{equation}
and the $O(\eps_{\rm BO}^2)$ term in Theorem~\ref{P13:thm:schrodinger}
is exactly zero.
\status{published}
\source{P13:thm:NA_vanishing}
\end{theorem}
\status{published}
\source{P13:thm:NA_vanishing}

\begin{corollary}[Improved error bound for the Schr\"odinger equation]
\label{P13:cor:improved_BO}
The Schr\"odinger equation~\eqref{P13:eq:schrodinger_derived} of
Theorem~\ref{P13:thm:schrodinger} holds with the improved error:
\begin{equation}
  \boxed{i\hbar\,\frac{\partial\psi}{\partial t}
  \;=\;
  \left[-\frac{\hbar^2}{2m_e}\nabla^2 + V(\mathbf{R})\right]\psi
  \;+\; O(\eps_{\rm cond}),}
  \label{P13:eq:schrodinger_improved}
\end{equation}
where $\eps_{\rm cond}\sim10^{-27}$ is the condensate anisotropy
from Theorem~\ref{P13:thm:modified_H}.
The formal $O(\eps_{\rm BO}^2)$ term vanishes exactly
(Theorem~\ref{P13:thm:NA_vanishing}), improving the error bound by
a factor of
\begin{equation}
  \boxed{\frac{\eps_{\rm BO}^2}{\eps_{\rm cond}}
  \;=\; \frac{[\aem^2/(2\vp)]^2}{\eps_{\rm cond}}
  \;\approx\; \frac{2.7\times10^{-10}}{2.2\times10^{-27}}
  \;\approx\; 1.2\times10^{17}.}
  \label{P13:eq:improvement_factor}
\end{equation}
The Schr\"odinger equation is therefore exact to one part in $10^{27}$,
not merely one part in $10^{10}$.
\status{published}
\source{P13:cor:improved_BO}
\end{corollary}
\status{published}
\source{P13:cor:improved_BO}

\begin{remark}[Why the condensate BO is exact and molecular BO is not]
\label{P13:rem:BO_comparison}
The contrast between the condensate and molecular cases is sharp:
\begin{enumerate}[noitemsep]
\item \emph{Molecule} (e.g.\ H$_2$):
  Electronic wavefunction $\Phi_0(r;R)$ depends on nuclear coordinate $R$.
  Derivative coupling $d_{0k}\neq0$ ($\sim\sqrt{m_e/M}$ per matrix element).
  Non-adiabatic correction $\sim\eps_{\rm BO}\sim m_e/M$.
\item \emph{Condensate gradient-knot} (this paper):
  Internal states in $\mathcal{H}_{\rm internal}$ are $R$-independent by the
  translation fibration~\eqref{P13:eq:fibration}.
  Derivative coupling $d_{0k}=0$ \emph{exactly} for all $k$.
  Non-adiabatic correction $=0$; residual error $=O(\eps_{\rm cond})$.
\end{enumerate}
The condensate BO is \emph{exact} in a sense that ordinary molecular BO is not:
the internal and external degrees of freedom are structurally decoupled
by the configuration-space fibration, not merely approximately so by a large mass ratio.
\status{published}
\source{P13:rem:BO_comparison}
\end{remark}
\status{published}
\source{P13:rem:BO_comparison}

\begin{remark}[Condensate phonons and the electron]
\label{P13:rem:phonons}
The condensate background has phonon excitations at
$\Delta_{\rm phonon}\sim\vp^2\Lcond\approx3\times10^{-4}\,\mathrm{eV}$,
which are \emph{below} the electron kinetic energy $\ERy\approx13.6\,\mathrm{eV}$.
The electron can in principle excite condensate phonons as it orbits.
However, the coupling of the electron gradient-knot to condensate
phonons is suppressed by the same chameleon factor
$\eps_{\rm cond}\sim10^{-27}$:
the atom's EM field density overwhelms the condensate background density,
screening all condensate interactions.
Net phonon emission rate: $\Gamma_{\rm phonon}\sim\eps_{\rm cond}^2\times
\text{(phase space)}\ll10^{-50}\,\mathrm{s}^{-1}$.
The electron's quantum mechanics is entirely standard at atomic energies.
\status{published}
\source{P13:rem:phonons}
\end{remark}
\status{published}
\source{P13:rem:phonons}

%════════════════════════════════════════════════════════════════════
% Supporting equations (referenced above but not independently boxed
% as named results)
%════════════════════════════════════════════════════════════════════

\begin{equation}
  E_{\rm cm}
  \;=\; \int\pi(x)\,\partial_0\rho(x)\,d^3x - \mathcal{L}
  \;\approx\; \tfrac{1}{2}M_{\rm eff}\,\mathbf{v}^2,
  \label{P13:eq:E_cm}
\end{equation}

\begin{equation}
  \Hcond
  \;=\; -\frac{\hbar^2}{2m_e}
        \,\frac{\epsr}{r^2}\,L^2,
  \label{P13:eq:H_cond}
\end{equation}

\begin{equation}
  \hat{H}_{\rm cond}
  \;=\; \hat{H}_{\rm cm}(\mathbf{R})
  \;+\; \hat{H}_{\rm internal}
  \;+\; \hat{V}(\mathbf{R}),
  \label{P13:eq:H_decomp}
\end{equation}

\begin{equation}
  \mathcal{H}
  \;=\; L^2\!\left(\mathcal{C}_Q^{\rm red},d\mu\right)
  \;\cong\;
  L^2\!\left(\mathbb{R}^3,d^3R\right)\otimes\mathcal{H}_{\rm internal},
  \label{P13:eq:Hilbert_factorisation}
\end{equation}

\begin{equation}
  \boxed{\epsr
  \;\equiv\; \frac{1}{\vp^6(1+\bst\rho(r))}}
  \label{P13:eq:epsilon_def}
\end{equation}

\begin{equation}
  \epsr
  \;\approx\; \frac{1}{\vp^6\,\bst\rho_{\rm atom}(r)}
  \;\approx\; \frac{2\Lcond^4}{\vp^6\,\aem^6 m_e^4}
        \left(\frac{r}{a_0}\right)^4,
  \label{P13:eq:eps_chameleon}
\end{equation}

\begin{equation}
  \mathcal{C}_Q^{\rm red} \;=\; \mathbb{R}^3 \times \mathcal{C}_{\rm internal},
  \label{P13:eq:fibration}
\end{equation}

\begin{equation}
  |\nabla\psi|^2
  \;=\; \underbrace{\left|\frac{\partial\psi}{\partial r}\right|^2}_{\text{radial}}
  \;+\; \underbrace{\frac{1}{r^2}|L\psi|^2}_{\text{angular}},
  \label{P13:eq:gradient_decomp}
\end{equation}

\begin{equation}
  \boxed{\Delta_{\rm internal}
  \;=\; m_e(\vp^{-38}-\vp^{-40})\big/\vp^{-40}
  \;=\; m_e(\vp^2-1)
  \;=\; \vp\, m_e
  \;\approx\; 827\,\mathrm{keV},}
  \label{P13:eq:internal_gap}
\end{equation}

\begin{equation}
  \boxed{\eps_{\rm BO}
  \;=\; \frac{\ERy}{\vp\, m_e}
  \;=\; \frac{\aem^2}{2\vp}
  \;\approx\; 1.6\times10^{-5}.}
  \label{P13:eq:eps_BO}
\end{equation}

\begin{equation}
  \hat{P}_i\,\psi(\mathbf{R})\otimes|0_{\rm internal}\rangle
  \;=\; \bigl(-i\hbar\,\partial/\partial R_i\,\psi(\mathbf{R})\bigr)
        \otimes|0_{\rm internal}\rangle.
  \label{P13:eq:momentum_on_sector_dup}
\end{equation}

\begin{equation}
  \mathcal{P}_i(x)
  \;=\; \pi(x)\,\partial_i\rho(x),
  \qquad
  \pi(x) \;=\; \frac{\partial\mathcal{L}}{\partial(\partial_0\rho)},
  \label{P13:eq:noether_momentum}
\end{equation}

\begin{equation}
  \mathcal{L} \;=\; \frac{|\nabla\rho|^2}{\vp^6(1+\bst|\nabla\rho|^2)}
  \label{P13:eq:parent_L}
\end{equation}

\begin{equation}
  \bst\rho_{\rm atom}(r)
  \;\sim\;
  \frac{u_{\rm EM}(r)}{\Lcond^4}
  \;=\; \frac{\aem^6 m_e^4}{2\Lcond^4}
        \left(\frac{a_0}{r}\right)^4.
  \label{P13:eq:rho_atom}
\end{equation}

\begin{equation}
  P_i \;=\; \int \pi(x)\,\partial_i\rho(x)\,d^3x.
  \label{P13:eq:total_momentum}
\end{equation}

\begin{equation}
  T_{\mathbf{a}}[\rho](x) = \rho(x-\mathbf{a}),
  \label{P13:eq:translation_action}
\end{equation}

\begin{equation}
  u_{\rm EM}(r)
  \;=\; \frac{\aem^2}{2r^4}
  \;=\; \frac{\aem^6 m_e^4}{2}\cdot\left(\frac{a_0}{r}\right)^4,
  \label{P13:eq:uEM}
\end{equation}
